% GENERATED FILE. Edit canonical lessons, then run npm run generate:books.
% canonical-source-hash: fnv1a64:838c6241dfbcc7f5
% canonical-lessons: HI-C556-sandeh, HI-C556-yaqinan, HI-C556-hargiz, HI-C556-vaham, HI-C556-niscay

\chapter{How Sure Am I?}
\label{ch:hi-how-sure-am-i}
\hlchaptermodality{556}

\begin{chapteropening}
\textbf{By the end of this chapter:} \emph{I can say how sure or unsure I am.}

Complete the last lesson of chapter 556: I can say how sure or unsure I am.
\end{chapteropening}

\section[\texorpdfstring{\hi{संदेह} (sandeh)}{संदेह (sandeh)}]{\hi{संदेह} (sandeh) — doubt (formal)}

\label{lesson:HI-C556-sandeh}



\begin{quote}
Before the new one: say the Hindi for a fall, then the Hindi for a downfall.
\end{quote}

\subsection*{You'll want to know: \hi{संदेह}}
\textbf{\hi{संदेह}} — \emph{sandeh} — \textquotedblleft{}doubt (formal)\textquotedblright{}.

\begin{quote}
\hi{उसकी} \hi{बात} \hi{पर} \hi{मुझे} \hi{कोई} \hi{संदेह} \hi{नहीं} \hi{है।}
\end{quote}

\begin{quote}
\emph{Uskī bāt par mujhe koī sandeh nahīṃ hai.}
\end{quote}

\begin{quote}
I have no doubt at all about what he says.
\end{quote}

Masculine.

A formal word; everyday speech says \hi{शक}.

\begin{cousinweb}
\textbf{\hi{संदेह}} (\emph{sandeh}) is a tatsama, taken whole from Sanskrit \emph{saṃdeha}, \textquotedblleft{}doubt, uncertainty\textquotedblright{} (Monier-Williams).
\end{cousinweb}

\begin{grammarlens}[title={What you've built}]
One more word for saying what happens if.
\end{grammarlens}

\subsection*{Your turn}
\begin{itemize}
\raggedright
  \item \emph{Say it:} \emph{sandeh}
  \item \emph{Say it:} \emph{sandeh}, once more
  \item \emph{Say it:} \emph{Uskī bāt par mujhe koī sandeh nahīṃ hai.} — I have no doubt at all about what he says.
  \item \emph{Recall:} say the Hindi for a fall, then the Hindi for a downfall, then say \emph{sandeh} again
\end{itemize}

\begin{tcolorbox}[breakable,colback=teal!4,colframe=teal!35!black,title={Before you move on}]
What does \hi{संदेह} mean? (\textquotedblleft{}Doubt (formal)\textquotedblright{}.) Say it once more.
\end{tcolorbox}
\section[\texorpdfstring{\hi{यक़ीनन} (yaqīnan)}{यक़ीनन (yaqīnan)}]{\hi{यक़ीनन} (yaqīnan) — certainly, surely}

\label{lesson:HI-C556-yaqinan}



\begin{quote}
Before the new one: say the Hindi for a downfall, then the Hindi for doubt (formal).
\end{quote}

\subsection*{You'll want to know: \hi{यक़ीनन}}
\textbf{\hi{यक़ीनन}} — \emph{yaqīnan} — \textquotedblleft{}certainly, surely\textquotedblright{}.

\begin{quote}
\hi{यक़ीनन} \hi{वह} \hi{कल} \hi{आएगा।}
\end{quote}

\begin{quote}
\emph{Yaqīnan vah kal āegā.}
\end{quote}

\begin{quote}
He will certainly come tomorrow.
\end{quote}

\begin{cousinweb}
\textbf{\hi{यक़ीनन}} (\emph{yaqīnan}) is usually traced to Arabic \emph{yaqīnan}, \textquotedblleft{}certainly\textquotedblright{}, through Persian.
\end{cousinweb}

\begin{grammarlens}[title={What you've built}]
One more word for saying what happens if.
\end{grammarlens}

\subsection*{Your turn}
\begin{itemize}
\raggedright
  \item \emph{Say it:} \emph{yaqīnan}
  \item \emph{Say it:} \emph{yaqīnan}, once more
  \item \emph{Say it:} \emph{Yaqīnan vah kal āegā.} — He will certainly come tomorrow.
  \item \emph{Recall:} say the Hindi for a downfall, then the Hindi for doubt (formal), then say \emph{yaqīnan} again
\end{itemize}

\begin{tcolorbox}[breakable,colback=teal!4,colframe=teal!35!black,title={Before you move on}]
What does \hi{यक़ीनन} mean? (\textquotedblleft{}Certainly, surely\textquotedblright{}.) Say it once more.
\end{tcolorbox}
\section[\texorpdfstring{\hi{हरगिज़} (hargiz)}{हरगिज़ (hargiz)}]{\hi{हरगिज़} (hargiz) — ever; (with \hi{नहीं}) never ever}

\label{lesson:HI-C556-hargiz}



\begin{quote}
Before the new one: say the Hindi for doubt (formal), then the Hindi for certainly, surely.
\end{quote}

\subsection*{You'll want to know: \hi{हरगिज़}}
\textbf{\hi{हरगिज़}} — \emph{hargiz} — \textquotedblleft{}ever; (with \hi{नहीं}) never ever\textquotedblright{}.

\begin{quote}
\hi{मैं} \hi{हरगिज़} \hi{झूठ} \hi{नहीं} \hi{बोलूँगा।}
\end{quote}

\begin{quote}
\emph{Maiṁ hargiz jhūṭh nahīṃ bolūṅgā.}
\end{quote}

\begin{quote}
I will never, ever lie.
\end{quote}

\begin{cousinweb}
\textbf{\hi{हरगिज़}} (\emph{hargiz}) is usually traced to Persian \emph{hargiz}, \textquotedblleft{}ever\textquotedblright{}.
\end{cousinweb}

\begin{grammarlens}[title={What you've built}]
One more word for saying what happens if.
\end{grammarlens}

\subsection*{Your turn}
\begin{itemize}
\raggedright
  \item \emph{Say it:} \emph{hargiz}
  \item \emph{Say it:} \emph{hargiz}, once more
  \item \emph{Say it:} \emph{Maiṁ hargiz jhūṭh nahīṃ bolūṅgā.} — I will never, ever lie.
  \item \emph{Recall:} say the Hindi for doubt (formal), then the Hindi for certainly, surely, then say \emph{hargiz} again
\end{itemize}

\begin{tcolorbox}[breakable,colback=teal!4,colframe=teal!35!black,title={Before you move on}]
What does \hi{हरगिज़} mean? (\textquotedblleft{}Ever; (with \hi{नहीं}) never ever\textquotedblright{}.) Say it once more.
\end{tcolorbox}
\section[\texorpdfstring{\hi{वहम} (vaham)}{वहम (vaham)}]{\hi{वहम} (vaham) — a groundless fear; a fancy}

\label{lesson:HI-C556-vaham}



\begin{quote}
Before the new one: say the Hindi for certainly, surely, then the Hindi for ever.
\end{quote}

\subsection*{You'll want to know: \hi{वहम}}
\textbf{\hi{वहम}} — \emph{vaham} — \textquotedblleft{}a groundless fear; a fancy\textquotedblright{}.

\begin{quote}
\hi{यह} \hi{तुम्हारा} \hi{वहम} \hi{है}, \hi{कोई} \hi{चोर} \hi{नहीं} \hi{है।}
\end{quote}

\begin{quote}
\emph{Yah tumhārā vaham hai, koī cor nahīṃ hai.}
\end{quote}

\begin{quote}
That is only in your head; there is no thief.
\end{quote}

Masculine.

\begin{cousinweb}
\textbf{\hi{वहम}} (\emph{vaham}) is usually traced to Arabic \emph{wahm}, \textquotedblleft{}a fancy, a supposition\textquotedblright{}, through Persian.
\end{cousinweb}

\begin{grammarlens}[title={What you've built}]
One more word for saying what happens if.
\end{grammarlens}

\subsection*{Your turn}
\begin{itemize}
\raggedright
  \item \emph{Say it:} \emph{vaham}
  \item \emph{Say it:} \emph{vaham}, once more
  \item \emph{Say it:} \emph{Yah tumhārā vaham hai, koī cor nahīṃ hai.} — That is only in your head; there is no thief.
  \item \emph{Recall:} say the Hindi for certainly, surely, then the Hindi for ever, then say \emph{vaham} again
\end{itemize}

\begin{tcolorbox}[breakable,colback=teal!4,colframe=teal!35!black,title={Before you move on}]
What does \hi{वहम} mean? (\textquotedblleft{}A groundless fear; a fancy\textquotedblright{}.) Say it once more.
\end{tcolorbox}
\section[\texorpdfstring{\hi{निश्चय} (niścay)}{निश्चय (niścay)}]{\hi{निश्चय} (niścay) — a firm decision; a resolve}

\label{lesson:HI-C556-niscay}



\begin{quote}
Before the new one: say the Hindi for ever, then the Hindi for a groundless fear.
\end{quote}

\subsection*{You'll want to know: \hi{निश्चय}}
\textbf{\hi{निश्चय}} — \emph{niścay} — \textquotedblleft{}a firm decision; a resolve\textquotedblright{}.

\begin{quote}
\hi{मैंने} \hi{रोज़} \hi{पढ़ने} \hi{का} \hi{निश्चय} \hi{किया।}
\end{quote}

\begin{quote}
\emph{Maiṁne roz paṛhne kā niścay kiyā.}
\end{quote}

\begin{quote}
I decided firmly to study every day.
\end{quote}

Masculine.

\begin{cousinweb}
\textbf{\hi{निश्चय}} (\emph{niścay}) is a tatsama, taken whole from Sanskrit \emph{niścaya}, \textquotedblleft{}certainty; resolve\textquotedblright{} (Monier-Williams).
\end{cousinweb}

\begin{grammarlens}[title={What you've built}]
One more word for saying what happens if.
\end{grammarlens}

\subsection*{Your turn}
\begin{itemize}
\raggedright
  \item \emph{Say it:} \emph{niścay}
  \item \emph{Say it:} \emph{niścay}, once more
  \item \emph{Say it:} \emph{Maiṁne roz paṛhne kā niścay kiyā.} — I decided firmly to study every day.
  \item \emph{Recall:} say the Hindi for ever, then the Hindi for a groundless fear, then say \emph{niścay} again
\end{itemize}

\begin{tcolorbox}[breakable,colback=teal!4,colframe=teal!35!black,title={Before you move on}]
What does \hi{निश्चय} mean? (\textquotedblleft{}A firm decision; a resolve\textquotedblright{}.) Say it once more.
\end{tcolorbox}
